\documentclass[conference]{IEEEtran}

% ── Paquetes ──────────────────────────────────────────────────────
\usepackage[utf8]{inputenc}
\usepackage[spanish,es-nodecimaldot]{babel}
\usepackage{graphicx}
\usepackage{amsmath,amssymb}
\usepackage{booktabs}
\usepackage{tabularx}
\usepackage[hidelinks]{hyperref}
\usepackage{float}
\usepackage{xcolor}
\usepackage{listings}
\usepackage{cite}
\usepackage{url}

% ── Estilo de listados de código ─────────────────────────────────
\lstset{
  basicstyle=\ttfamily\scriptsize,
  keywordstyle=\color{blue!70!black}\bfseries,
  commentstyle=\color{green!40!black}\itshape,
  stringstyle=\color{red!60!black},
  language=C++,
  showstringspaces=false,
  numbers=left,
  numberstyle=\tiny\color{gray},
  frame=single,
  breaklines=true,
  tabsize=2,
  captionpos=b
}

% ── Metadatos ────────────────────────────────────────────────────
\title{Interpolación Geométrica de Contornos de Gliomas entre
Slices de Resonancia Magnética}

\author{\IEEEauthorblockN{Abel Albuez, Santiago Gil Gallego y
Jesús Alberto}
\IEEEauthorblockA{Maestría en Ingeniería de Sistemas y Computación\\
Pontificia Universidad Javeriana, Bogotá D.C., Colombia\\
Profesor: Leonardo Flórez-Valencia \quad
Geometría Computacional \quad 2026\\
Repositorio: \url{https://github.com/AbelAlbuez/computational-geometry-cgal}}}

\begin{document}

\maketitle

% ── Abstract ─────────────────────────────────────────────────────
\begin{abstract}
Se presenta un módulo de interpolación geométrica entre
contornos 2D del tumor activo (\textit{Enhancing Tumor}, ET)
extraídos de cortes axiales consecutivos de resonancia
magnética cerebral. Sobre el dataset público
\textit{BraTS 2024 GLI}, los contornos se obtienen por slice
mediante segmentación experta y se almacenan como polilíneas
cerradas en formato \texttt{.obj}. El módulo, implementado en
C++17/20 sobre CGAL, encadena (i) un \emph{resampling} por
longitud de arco que iguala el número de vértices de ambos
contornos, (ii) una interpolación lineal vértice a vértice y
(iii) una etapa de detección y resolución de
auto-intersecciones basada en el algoritmo de
Bentley--Ottmann provisto por la librería \texttt{pujCGAL}.
Este informe describe el módulo desarrollado por el
integrante \textbf{Abel Albuez}, validado sobre el caso
\textit{BraTS-GLI-00008-100}, slices~70 y~71, y se enmarca
en una arquitectura mayor cuyos demás componentes son
responsabilidad de los otros integrantes del equipo.
\end{abstract}

\begin{IEEEkeywords}
interpolación de contornos, geometría computacional, CGAL,
Bentley--Ottmann, BraTS, resonancia magnética, gliomas
\end{IEEEkeywords}

% =====================================================================
\section{Introducción}
% =====================================================================
La reconstrucción tridimensional de tumores cerebrales a partir
de secuencias de resonancia magnética (MRI) es una tarea
relevante para el diagnóstico, la planeación quirúrgica y el
seguimiento clínico de pacientes con gliomas. Una estrategia
extendida consiste en segmentar slice por slice la región de
interés y luego \emph{interpolar} los contornos 2D obtenidos
para reconstruir una superficie cerrada en 3D.

Este trabajo utiliza el dataset \textit{BraTS 2024
GLI}~\cite{baid2021brats} (Brain Tumor Segmentation Challenge
--- Adult Glioma), que provee volúmenes de MRI con
segmentaciones expertas. De cada volumen se extrae el contorno
del \textit{Enhancing Tumor} (label ET~$=3$) por slice axial
mediante \texttt{scikit-image}, obteniendo polígonos cerrados
almacenados en formato \texttt{.obj}.

La motivación última proviene de la línea de investigación del
autor sobre segmentación automática de gliomas: una vez
segmentado el tumor en slices 2D, se requiere un módulo
geométrico que produzca slices intermedios consistentes
---sin auto-intersecciones--- para alimentar etapas
posteriores de reconstrucción volumétrica.

% =====================================================================
\section{Marco Teórico}
\label{sec:marco}
% =====================================================================

\subsection{Interpolación entre slices paralelos}
Barequet y Sharir~\cite{barequet1996} formalizan el problema
de interpolación lineal por tramos entre cortes poligonales
paralelos como un problema de \emph{correspondencia} entre
vértices: dadas dos poligonales $A$ y $B$ en planos paralelos,
construir una superficie triangulada que las una. Su trabajo
establece el marco bajo el cual opera el módulo aquí descrito,
restringido al caso $|A| = |B|$ (igualando los vértices mediante
\textit{resampling}) y a la interpolación lineal directa de
vértices correspondientes.

\subsection{Interpolación en imágenes médicas}
Meijster y Roerdink~\cite{meijster1999} comparan métodos de
interpolación (vecino más cercano, lineal, cúbica, basada en
forma) en el contexto de procesamiento de imágenes médicas.
Concluyen que la interpolación lineal por vértices ofrece una
relación calidad/complejidad adecuada para contornos densamente
muestreados, justificando su uso como \textit{baseline} en este
módulo antes de considerar variantes basadas en forma o en
transformadas de distancia.

\subsection{Polígonos auto-intersecantes}
Souviron~\cite{souviron2013} presenta un esquema de detección
y corrección de auto-intersecciones en polígonos cerrados con
memoria mínima. En este trabajo se adopta una estrategia
equivalente en el espíritu pero apoyada en CGAL: se detectan
los cruces con el barrido de Bentley--Ottmann y se descompone
el polígono en sub-lazos, reteniendo el de mayor área. El
algoritmo de Bentley--Ottmann en sí está implementado dentro
de la librería \texttt{pujCGAL} provista por el profesor del
curso.

% =====================================================================
\section{Metodología --- Módulo de Abel}
\label{sec:metodologia}
% =====================================================================
El módulo está estructurado como una cadena de clases
independientes sobre el kernel
\texttt{CGAL::Exact\_predicates\_inexact\_constructions\_kernel}
y el tipo \texttt{TContour}~$=$
\texttt{std::vector<Point\_2>}. Las responsabilidades se
descomponen en cuatro etapas, cada una encapsulada en una
clase con un único método estático.

% ---------------------------------------------------------------------
\subsection{Resampling por longitud de arco
            (\texttt{ContourResampler})}
% ---------------------------------------------------------------------
Dado un contorno cerrado $C = \{p_0, p_1, \dots, p_{m-1}\}$,
se calcula la longitud de arco acumulada
\begin{equation}
  s_0 = 0,
  \qquad
  s_{i+1} = s_i + \lVert p_{(i+1) \bmod m} - p_i \rVert,
  \quad i = 0, \dots, m-1,
\end{equation}
con $L = s_m$ la longitud total del contorno cerrado.

Para producir un contorno con exactamente $n$ vértices
uniformemente distribuidos por longitud de arco se evalúa,
para cada $k \in \{0,\dots,n-1\}$, el parámetro
$s_k^{\star} = k \, L / n$, se localiza el segmento
$[p_i, p_{i+1}]$ tal que $s_i \leq s_k^{\star} \leq s_{i+1}$,
y se interpola
\begin{equation}
  u = \frac{s_k^{\star} - s_i}{s_{i+1} - s_i},
  \qquad
  q_k = (1 - u)\, p_i + u \, p_{(i+1) \bmod m}.
\end{equation}
Esta etapa garantiza la precondición $|A| = |B|$ requerida por
la interpolación lineal.

% ---------------------------------------------------------------------
\subsection{Interpolación lineal
            (\texttt{LinearInterpolator})}
% ---------------------------------------------------------------------
Sean $A = \{a_0, \dots, a_{n-1}\}$ y
$B = \{b_0, \dots, b_{n-1}\}$ dos contornos con el mismo
número de vértices (tras el \textit{resampling}). El contorno
interpolado $P(t)$ para $t \in [0,1]$ se define
componente a componente:
\begin{equation}
  P_i(t) = (1 - t)\, a_i + t \, b_i,
  \qquad i = 0, 1, \dots, n-1.
\end{equation}
En particular, $t = 0{,}5$ produce el slice intermedio entre
$A$ y $B$.

\begin{lstlisting}[caption={Núcleo de la interpolación lineal.}]
for( std::size_t i = 0; i < n; ++i )
{
  const double ax = CGAL::to_double( A[i].x() );
  const double ay = CGAL::to_double( A[i].y() );
  const double bx = CGAL::to_double( B[i].x() );
  const double by = CGAL::to_double( B[i].y() );
  out.emplace_back( (1.0 - t)*ax + t*bx,
                    (1.0 - t)*ay + t*by );
}
\end{lstlisting}

% ---------------------------------------------------------------------
\subsection{Detección de auto-intersecciones
            (\texttt{SelfIntersectionResolver})}
% ---------------------------------------------------------------------
El contorno interpolado se convierte primero en un vector de
segmentos CGAL conectando vértices consecutivos en orden
circular:
\begin{equation}
  S_i = \overline{P_i \, P_{(i+1) \bmod n}},
  \qquad i = 0,\dots,n-1.
\end{equation}
A continuación se invoca el algoritmo de barrido
Bentley--Ottmann~\cite{souviron2013} provisto por
\texttt{pujCGAL}:

\begin{lstlisting}[caption={Invocación del barrido
Bentley--Ottmann de \texttt{pujCGAL}.}]
std::vector< TPoint > inters;
pujCGAL::SegmentsIntersection::BentleyOttmann(
  segs.begin(), segs.end(),
  std::back_inserter( inters )
);
\end{lstlisting}

Se filtran las intersecciones que coinciden con un vértice
original del contorno (extremos compartidos entre aristas
consecutivas, que no representan cruces reales). Si el
conjunto resultante es vacío, el contorno se retorna
inalterado.

% ---------------------------------------------------------------------
\subsection{Resolución de cruces}
% ---------------------------------------------------------------------
Cuando existen cruces, cada punto de intersección $p_j$ se
asocia a las dos aristas que lo contienen y se calcula el
parámetro $t \in [0,1]$ de su posición a lo largo de cada
arista. Esto permite construir una \emph{traversal subdividida}:
una lista ordenada que recorre el contorno insertando los
puntos de cruce en su posición correspondiente.

Sobre esta traversal se aplica un esquema clásico de
extracción de lazos mediante pila: cuando un punto de
intersección aparece por segunda vez se cierra un sub-lazo
con todos los nodos apilados desde su primera aparición. Cada
sub-lazo $L_k$ tiene un área absoluta dada por la fórmula del
polígono (\emph{shoelace}):
\begin{equation}
  \mathrm{Area}(L_k) =
  \frac{1}{2}
  \left|
    \sum_{i=0}^{|L_k|-1}
    (x_i \, y_{i+1} - x_{i+1} \, y_i)
  \right|.
\end{equation}
El módulo retiene el sub-lazo $L_{k^{\star}}$ con
$k^{\star} = \arg\max_k \mathrm{Area}(L_k)$, es decir, el lazo
exterior de mayor extensión, descartando lazos parásitos
típicos cuando los contornos fuente difieren mucho en forma.

% =====================================================================
\section{Resultados}
\label{sec:resultados}
% =====================================================================

\subsection{Configuración experimental}
\label{sec:resultados:config}
La evaluación se realizó sobre el dataset
\textit{BraTS 2024 GLI} restringido al label
\textit{Enhancing Tumor} (ET, valor~$3$). Los contornos se
extrajeron por slice axial mediante \texttt{scikit-image} y
se serializaron en formato \texttt{.obj} 2D. Sobre este
corpus se ejecutaron \textbf{4 pares representativos} de
slices con parámetro de interpolación fijo $t = 0{,}5$,
cada uno con una intención distinta:
\begin{enumerate}
  \item slices consecutivos del mismo caso (variación
        morfológica mínima);
  \item slices separados por 6 cortes del mismo caso
        (variación moderada);
  \item par extremo (primer y último slice con ET) dentro de
        un mismo caso (variación máxima intra-caso);
  \item par cruzado entre dos casos distintos (primer slice
        del caso $X$ contra primer slice del caso $Y$).
\end{enumerate}
Sobre la versión inicial del interpolador se incorporaron
\textbf{tres mejoras geométricas} de pre-alineación
(orientación, centroide y rotación cíclica óptima del punto
de inicio), descritas en la Sección~\ref{sec:resultados:mejoras}.
Cada ejecución produce un \texttt{.obj} 3D y una
visualización PNG con los dos contornos fuente y el contorno
interpolado superpuestos.

\subsection{Mejoras geométricas implementadas}
\label{sec:resultados:mejoras}

\subsubsection{Normalización de orientación (CCW)}
\label{sec:resultados:mejoras:ccw}
La interpolación lineal vértice a vértice es
\emph{sensible al sentido de recorrido}: si $A$ está en
sentido CCW y $B$ en CW, la correspondencia $A_i
\leftrightarrow B_i$ recorre los contornos en direcciones
opuestas y produce un polígono retorcido. Para evitarlo se
calcula el área con signo de cada contorno mediante la
fórmula \emph{shoelace}:
\begin{equation}
  \mathrm{area}(C) =
  \tfrac{1}{2}\!\sum_{i=0}^{n-1}\!
  \bigl( x_i\, y_{i+1} - x_{i+1}\, y_i \bigr),
\end{equation}
y si $\mathrm{area}(C) < 0$ se invierte el orden de los
vértices, garantizando que ambos contornos sean CCW antes
de interpolar.

\subsubsection{Alineación por centroide}
\label{sec:resultados:mejoras:centroide}
La interpolación lineal implícitamente promedia las
posiciones absolutas. Para que la geometría intermedia no
quede distorsionada por una traslación entre $A$ y $B$, se
calculan los centroides
\begin{equation}
  \bar{c}_A = \frac{1}{|A|}\sum_{i} A_i,
  \qquad
  \bar{c}_B = \frac{1}{|B|}\sum_{i} B_i,
\end{equation}
se trasladan ambos contornos al origen, se interpola en el
sistema centrado y el contorno resultante se traslada al
centroide promedio $\bar{c}_M = \tfrac{1}{2}(\bar{c}_A +
\bar{c}_B)$.

\subsubsection{Rotación óptima del punto de inicio}
\label{sec:resultados:mejoras:rotacion}
Aún con $|A| = |B|$ tras el \textit{resampling}, el primer
vértice de cada contorno es arbitrario. Una correspondencia
$A_i \leftrightarrow B_i$ desalineada estira
innecesariamente las aristas del polígono interpolado. Se
busca el desplazamiento cíclico $k^{\star}$ que minimiza la
suma de distancias al cuadrado:
\begin{equation}
  k^{\star} \;=\; \arg\min_{k \in \{0,\dots,n-1\}}\;
  \sum_{i=0}^{n-1}
  \bigl\lVert A_i - B_{(i+k) \bmod n} \bigr\rVert^{2},
\end{equation}
y se aplica $B \leftarrow \mathrm{rotate}(B, k^{\star})$
antes de interpolar.

\subsection{Tabla de resultados}
\label{sec:resultados:tabla}
La Tabla~\ref{tab:resultados} reporta los conteos de
vértices, la rotación óptima detectada $k^{\star}$ y el
estado de auto-intersección para los cuatro pares. La
columna \emph{Resam.} indica el número común de vértices
tras el \textit{resampling}, igual a $\max(|A|,|B|)$.

\begin{table}[!t]
\centering
\caption{Resultados de los 4 pares representativos
($t = 0{,}5$) tras aplicar las mejoras de pre-alineación.}
\label{tab:resultados}
\renewcommand{\arraystretch}{1.15}
\setlength{\tabcolsep}{4pt}
\footnotesize
\begin{tabular}{@{}llrrrrc@{}}
\toprule
\textbf{Caso} & \textbf{A$\rightarrow$B} &
\textbf{$|A|$} & \textbf{$|B|$} & \textbf{Resam.} &
\textbf{$k^{\star}$} & \textbf{Auto-int.} \\
\midrule
00008-100 (consecutivos) & 0070\,$\rightarrow$\,0071 &
  32 & 30 & 32 & 0 & No \\
00008-100 (separados 6)  & 0066\,$\rightarrow$\,0072 &
  14 & 24 & 24 & 1 & No \\
00008-101 (extremos)     & 0070\,$\rightarrow$\,0072 &
  26 & 16 & 26 & 2 & No \\
100 \emph{vs} 101 (cruzado) & 0066\,$\rightarrow$\,0070 &
  14 & 26 & 26 & 24 & No \\
\bottomrule
\end{tabular}
\end{table}

En los cuatro casos el barrido de Bentley--Ottmann no
reportó auto-intersecciones, por lo que el
\texttt{SelfIntersectionResolver} retornó el contorno
original sin modificación. Los cuatro contornos
interpolados resultantes son polígonos simples válidos. La
salida cruda del ejecutable para el primer par ilustra el
formato textual:

\begin{lstlisting}[language=bash, numbers=none,
                   caption={Salida de \texttt{contour\_interpolator}
                   sobre dos slices consecutivos.}]
Contour A: 32 vertices
Contour B: 30 vertices
Resampled to: 32 vertices
Optimal rotation of B: 0
Interpolated contour: 32 vertices
Self-intersections detected: no
Result written to output/caso1_consecutivos/contour_interpolated.obj
\end{lstlisting}

\subsection{Análisis por caso}
\label{sec:resultados:porcaso}

\textbf{Caso~1 — consecutivos (0070\,$\rightarrow$\,0071,
$k^{\star}=0$).} Resultado correcto: el contorno interpolado
queda visiblemente entre $A$ y $B$ y conserva la forma
general del tumor. La rotación óptima nula indica que las
parametrizaciones por longitud de arco ya estaban
naturalmente alineadas; persiste un artefacto leve en la
zona derecha por la diferencia de forma local entre los dos
slices.

\textbf{Caso~2 — separados 6 cortes
(0066\,$\rightarrow$\,0072, $k^{\star}=1$).} Los contornos
$A$ y $B$ ocupan zonas distintas del plano $xy$. La
alineación por centroide es la que más aporta aquí: sin
ella, el interpolado quedaría sesgado hacia el contorno
mayor. Con ella, el resultado se ubica en el centroide
promedio $\bar{c}_M$ y es geométricamente válido, aunque
sin correspondencia anatómica directa con un slice $z$
intermedio.

\textbf{Caso~3 — extremos del mismo caso
(0070\,$\rightarrow$\,0072, $k^{\star}=2$).} Mejor
resultado cualitativo. La rotación óptima no trivial
($k^{\star}=2$) y la alineación por centroide producen un
contorno intermedio que captura bien la transición de forma
entre el primer y el último slice con ET del caso. Es el
caso donde la mejora de rotación se hace más evidente al
comparar contra la versión sin pre-alineación.

\textbf{Caso~4 — cruzado entre casos
(0066\,$\rightarrow$\,0070, $k^{\star}=24$).} Limitación
intrínseca del método. La rotación elevada ($k^{\star}=24$
sobre $n=26$, equivalente a casi una vuelta completa)
revela que no existe una correspondencia natural entre los
vértices de contornos provenientes de casos distintos. La
salida es un polígono simple sin auto-intersecciones, pero
sin significado clínico: representa el promedio geométrico
entre dos tumores no comparables.

\subsection{Limitación documentada}
\label{sec:resultados:limitacion}
La interpolación lineal vértice a vértice asume
implícitamente que los contornos están alineados
espacialmente y representan instancias progresivas de una
misma estructura. Las mejoras de Sección~\ref{sec:resultados:mejoras}
robustecen esta hipótesis para slices del mismo caso, pero
no la sustituyen: cuando los contornos pertenecen a
pacientes distintos o están muy desplazados, el método
produce una salida geométricamente válida ---polígono
simple, sin auto-intersecciones--- pero clínicamente
irrelevante. Métodos basados en \emph{morphing}, distancia
geodésica o transformadas de distancia con signo resolverían
mejor este escenario y se consideran como extensión natural
del módulo.

\subsection{Visualizaciones}
\label{sec:resultados:figuras}
Las Figuras~\ref{fig:ref} y~\ref{fig:extremo_cruzado}
muestran los cuatro PNGs generados con la versión mejorada
del interpolador. En cada gráfico se dibujan los contornos
fuente $A$ (azul punteado) y $B$ (rojo punteado), y el
contorno interpolado a $t = 0{,}5$ (verde sólido con
marcadores).

\begin{figure}[H]
\centering
\includegraphics[width=0.46\columnwidth]{../output/caso1_consecutivos/caso1_consecutivos.png}
\includegraphics[width=0.46\columnwidth]{../output/caso2_separados/caso2_separados.png}
\caption{Caso \textit{BraTS-GLI-00008-100} con las mejoras
de pre-alineación activas. Izquierda: slices consecutivos
$0070 \rightarrow 0071$, $k^{\star}=0$. Derecha: slices
separados por 6 cortes, $0066 \rightarrow 0072$,
$k^{\star}=1$. Ambos resueltos sin auto-intersecciones.}
\label{fig:ref}
\end{figure}

\begin{figure}[H]
\centering
\includegraphics[width=0.46\columnwidth]{../output/caso3_extremos/caso3_extremos.png}
\includegraphics[width=0.46\columnwidth]{../output/caso4_cruzado/caso4_cruzado.png}
\caption{Izquierda: par extremo del caso
\textit{BraTS-GLI-00008-101}, primer ($0070$) contra último
slice con ET ($0072$), $k^{\star}=2$ (mejor caso
cualitativo). Derecha: par cruzado entre casos
\textit{BraTS-GLI-00008-100} ($0066$) y
\textit{BraTS-GLI-00008-101} ($0070$), $k^{\star}=24$
(limitación del método ante contornos no alineados
espacialmente). Ambos resueltos sin auto-intersecciones.}
\label{fig:extremo_cruzado}
\end{figure}

El archivo de salida se escribe en formato \texttt{.obj}
3D con $z = 0{,}0$ por vértice, lo que asegura compatibilidad
con visores estándar como ParaView, Online 3D Viewer e
ImageToSTL.

% =====================================================================
\section{Decisiones de diseño y problemas resueltos}
\label{sec:decisiones}
% =====================================================================

\subsection{Compatibilidad C++17 $\rightarrow$ C++20}
\label{sec:decisiones:cxx20}
Durante la compilación inicial el proyecto fallaba con el
error \texttt{'iter\_value\_t' is not a member of 'std'}. La
causa es que \texttt{pujCGAL/SegmentsIntersection.hxx} usa
\texttt{std::iter\_value\_t}, disponible únicamente desde
C++20. Para no modificar la librería del profesor se
agregó al \texttt{CMakeLists.txt} del módulo:
\begin{lstlisting}[caption={Activación de C++20 sólo para el
target del módulo.}]
target_compile_features(
  contour_interpolator PRIVATE cxx_std_20
);
\end{lstlisting}
Esto eleva el estándar exclusivamente para el ejecutable
\texttt{contour\_interpolator} y resuelve la compilación sin
tocar código externo.

\subsection{Formato \texttt{.obj} 2D \emph{vs} 3D}
\label{sec:decisiones:obj3d}
La primera versión del exportador escribía vértices en 2D
(\texttt{v x y}). Visores 3D como ParaView e ImageToSTL
rechazaban el archivo con errores del tipo
\emph{``Invalid vertex index''} o \emph{``Failed to import
model''}. La solución fue introducir una función local
\texttt{write\_obj\_3d} que agrega $z = 0{,}0$ a cada
vértice, produciendo el formato estándar \texttt{v x y 0.0}
compatible con cualquier visor 3D, sin necesidad de
transformaciones posteriores.

\subsection{Ausencia de auto-intersecciones en contornos
            BraTS}
\label{sec:decisiones:noselfint}
Durante el desarrollo se ejecutó un barrido amplio de
\textbf{94 pares} de slices (consecutivos, separados,
extremos y cruzados entre casos), sin detectar ninguna
auto-intersección. El análisis geométrico indica que los
contornos ET de BraTS son aproximadamente convexos y
morfológicamente suaves: el promedio ponderado de dos
contornos convexos alineados por longitud de arco tiende a
producir un contorno simple. Las tres mejoras de
pre-alineación descritas en
Sección~\ref{sec:resultados:mejoras} (normalización CCW,
alineación por centroide y rotación óptima del punto de
inicio) \emph{preservan} esta propiedad: ninguna introduce
nuevas aristas ni modifica la topología del contorno
interpolado, sólo reasignan correspondencias o aplican
rigid transforms (traslación, inversión, rotación cíclica).
En consecuencia, los 4 pares del experimento final
mantuvieron la validez topológica. El módulo
\texttt{SelfIntersectionResolver} está correctamente
implementado (detección con Bentley--Ottmann, partición en
sub-lazos y selección por área máxima) pero no se activa
con datos reales de BraTS. Para demostrar su funcionamiento
de forma controlada se planea un experimento con
correspondencia de vértices deliberadamente errónea
(\emph{modo stress}: rotación del array $B$ o rotación
geométrica previa), donde se garantiza la aparición de
cruces.

\subsection{Independencia del módulo respecto al equipo}
\label{sec:decisiones:independencia}
El módulo se diseñó para correr directamente sobre los
\texttt{.obj} extraídos de BraTS, sin depender del output de
los otros integrantes del equipo (Jesús: correspondencia de
vértices por ángulo polar; Santiago: reconstrucción 3D con
VTK). Esta decisión reduce el riesgo de integración y
permite validar el componente de Abel en cualquier momento
sin coordinación con el resto del equipo. La integración se
realiza a posteriori a través del tipo común
\texttt{TContour} y del formato \texttt{.obj}, ambos
estables y bien definidos.

% =====================================================================
\section{Conclusiones Parciales}
\label{sec:conclusiones}
% =====================================================================
\begin{itemize}
  \item El pipeline corre de extremo a extremo sobre datos
        reales del dataset \textit{BraTS 2024 GLI}, generando
        un contorno interpolado válido en formato
        \texttt{.obj} 3D que se renderiza en visores externos
        sin transformaciones adicionales.
  \item La librería \texttt{pujCGAL} provee el barrido de
        Bentley--Ottmann en una interfaz tipo \emph{iterator},
        lo cual evita reimplementar el algoritmo y permite
        concentrar el esfuerzo del módulo en la lógica de
        partición y selección de sub-lazos.
  \item Cada clase (\texttt{ContourResampler},
        \texttt{LinearInterpolator},
        \texttt{SelfIntersectionResolver}) tiene una
        responsabilidad única y opera sobre el tipo común
        \texttt{TContour}. Esto hace al módulo independiente
        de los componentes desarrollados por los otros
        integrantes del equipo y permite integrarlo sin
        acoplamiento adicional.
  \item Como trabajo futuro inmediato se contempla evaluar el
        módulo con pares de slices que produzcan
        auto-intersecciones reales ---típicamente en los
        cortes superiores e inferiores del tumor, donde la
        topología cambia bruscamente--- y comparar la
        selección por área máxima contra estrategias basadas
        en orientación o en distancia de Hausdorff entre
        sub-lazos candidatos.
\end{itemize}

% =====================================================================
% ── Referencias ──────────────────────────────────────────────────
% =====================================================================
\begin{thebibliography}{9}

\bibitem{barequet1996} G.~Barequet and M.~Sharir,
``Piecewise-linear interpolation between polygonal slices,''
\textit{Computer Vision and Image Understanding},
vol.~63, no.~2, pp.~251--272, 1996.

\bibitem{meijster1999} A.~Meijster and J.~B.~T.~M.~Roerdink,
``A survey of interpolation methods in medical image
processing,'' \textit{IEEE Transactions on Medical Imaging},
1999.

\bibitem{souviron2013} J.~Souviron, ``Correcting
self-intersecting polygons using minimal memory,''
\textit{arXiv:1305.4573}, 2013.

\bibitem{baid2021brats} U.~Baid \textit{et al.}, ``The RSNA-ASNR-MICCAI
BraTS 2021 benchmark on brain tumor segmentation and
radiogenomic classification,'' \textit{arXiv:2107.02314}, 2021.

\end{thebibliography}

\end{document}
